\documentclass[11pt,a4paper]{article}
\usepackage[margin=26mm,headheight=15pt]{geometry}
\usepackage[T1]{fontenc}
\usepackage[utf8]{inputenc}
\usepackage{libertinus}
\usepackage{amsmath,amssymb,amsthm,mathtools}
\usepackage{microtype,booktabs,array,longtable}
\usepackage{xcolor,enumitem,fancyhdr,titlesec,xurl}
\usepackage{aliascnt}
\usepackage{hyperref}
\usepackage[nameinlink,noabbrev]{cleveref}
\definecolor{ink}{RGB}{25,53,75}
\definecolor{muted}{RGB}{85,90,96}
\hypersetup{colorlinks=true,linkcolor=ink,citecolor=ink,urlcolor=ink,
 pdftitle={Fractional Cusps at the Atomic Phase},
 pdfauthor={Research draft prepared with ChatGPT for Vladimir Reshetnikov},
 pdfsubject={Sharp noninteger regularity and relaxation for digital-product pressure}}
\titleformat{\section}{\Large\bfseries\color{ink}}{\thesection}{.7em}{}
\titleformat{\subsection}{\large\bfseries}{\thesubsection}{.7em}{}
\pagestyle{fancy}\fancyhf{}
\fancyhead[L]{\small Fractional cusps at the atomic phase}
\fancyhead[R]{\small ProveIt research}
\fancyfoot[C]{\thepage}
\setlength{\emergencystretch}{3em}
\setlength{\parskip}{.2em}
\setlist{itemsep=.2em,topsep=.4em}
\allowdisplaybreaks
\numberwithin{equation}{section}
\newtheorem{theorem}{Theorem}[section]
\newaliascnt{lemma}{theorem}
\newtheorem{lemma}[lemma]{Lemma}
\aliascntresetthe{lemma}
\newaliascnt{proposition}{theorem}
\newtheorem{proposition}[proposition]{Proposition}
\aliascntresetthe{proposition}
\newaliascnt{corollary}{theorem}
\newtheorem{corollary}[corollary]{Corollary}
\aliascntresetthe{corollary}
\theoremstyle{definition}
\newaliascnt{definition}{theorem}
\newtheorem{definition}[definition]{Definition}
\aliascntresetthe{definition}
\newtheorem{question}{Research question}
\theoremstyle{remark}
\newaliascnt{remark}{theorem}
\newtheorem{remark}[remark]{Remark}
\aliascntresetthe{remark}
\newcommand{\TT}{\mathbb T}
\newcommand{\RR}{\mathbb R}
\newcommand{\CC}{\mathbb C}
\newcommand{\ZZ}{\mathbb Z}
\newcommand{\NN}{\mathbb N}
\newcommand{\ee}{\mathrm e}
\newcommand{\ii}{\mathrm i}
\newcommand{\dd}{\,\mathrm d}
\newcommand{\Lop}{\mathcal L}
\newcommand{\Proj}{\Pi}
\newcommand{\sgn}{\operatorname{sgn}}
\newcommand{\sinc}{\operatorname{sinc}}
\newcommand{\spec}{\operatorname{spec}}
\newcommand{\dist}{\operatorname{dist}}
\newcommand{\e}[1]{\ee^{2\pi\ii #1}}
\newcommand{\norm}[1]{\left\lVert#1\right\rVert}
\newcommand{\code}[1]{\texttt{\detokenize{#1}}}

\begin{document}
\begin{titlepage}
\thispagestyle{empty}
\setlength{\parindent}{0pt}
{\small\sffamily\color{ink} PROVEIT\quad /\quad ANALYSIS AND DYNAMICAL SYSTEMS}\par
\vspace{18mm}
{\Huge\bfseries\color{ink} Fractional Cusps\\[3mm] at the Atomic Phase\par}
\vspace{8mm}
{\Large Sharp noninteger regularity of Thue--Morse pressure\\[2mm]
and sharp relaxation for digital products in every base\par}
\vspace{11mm}
{\large Research draft prepared for Vladimir Reshetnikov\par}
{Prepared with ChatGPT\quad\textbullet\quad 29 September 2026\par}
\vspace{10mm}
\hrule\vspace{4mm}
\begin{abstract}
For the doubling map and the singular potential
$\psi_c(x)=\log\cos^2\!\pi(x-c)$, we determine the phase regularity of
$P(q\psi_c)$ at the atomic phase for every $q>1/2$. For every noninteger
$q$ in this range the pressure is locally $C^{\lceil2q\rceil-1}$, but its
$\lceil2q\rceil$th derivative at the atomic phase does not exist. Its leading
nonanalytic term is explicit:
\[
 P(q\psi_c)=2q\log\cos\pi c+
 2(2^{2q}-1)\zeta(2q)|c|^{2q}+o(|c|^{2q}).
\]
Thus analyticity at integer orders is exceptional, not evidence for
analyticity at all positive orders.

The proof constructs a base-independent, periodized-sinc eigenfunction and
establishes a sharp operator-norm relaxation law on every H\"older space.
This supplies the spectral gap directly, without assuming a general
Perron theorem for a potential with zeros. A positive-eigenfunction
perturbation argument, a pressure--integral comparison, and one exact
identity at the fixed point then give the singular expansion. The same
argument works in every integer base $b\ge2$, with cusp coefficient
$2(b^{2q}-1)\zeta(2q)$. At positive integer $q=m$, the would-be cusp cancels
exactly the $2m$th coefficient of the analytic fixed-branch contribution,
recovering and extending the repository's missing-coefficient phenomenon.

The article includes full proofs, a sharpness argument for the convergence
rate, finite-size moment asymptotics, numerical diagnostics, explicit
provenance and limitations, and a research agenda. The results are claims
proved in this research draft; they are not independently refereed or
proof-assistant verified.
\end{abstract}
\vfill
{\small\color{muted} Repository snapshot inspected:
\nolinkurl{afb2d1227d8bc5df3beffc1463e958db3544960b}.\\
The package contains the source, compiled PDF, executable checks, and their output.}
\end{titlepage}
\begingroup
\small
\setlength{\parskip}{0pt}
\tableofcontents
\endgroup
\newpage

\section{The question, the answer, and the scope}
\label{sec:intro}

\subsection{A published regularity question}
Gohlke, Kesseb\"ohmer, and Schindler prove order-two phase analyticity for
this pressure family and leave stronger regularity at general positive
orders open \cite[Theorem~2.7 and subsequent discussion]{GKS}.
Gohlke, Lamprinakis, and Schmeling subsequently prove continuity for all
nonnegative orders and explicitly repeat the stronger-regularity question
\cite[Theorems~2.2--2.3 and the following paragraph]{GLS}.
The latter paper uses $\log\sin^2\pi(x-a)$; its atomic phase $a=1/2$
corresponds to our phase $c=0$, by $a=c+1/2\pmod1$.

This article settles the local question at that exceptional phase for all
orders strictly above $1/2$. There is both a positive and a negative answer:
there are as many continuous derivatives as the order of vanishing of the
weight permits, but one further derivative fails at every noninteger
order. In particular, global real-analyticity at every positive order is
false. We do not claim to classify regularity at all other phases or at
orders at or below $1/2$.

\subsection{The connection to ProveIt}
The inspected repository contains a substantial Fabius/Rvachev and
Thue--Morse analysis program \cite{Repo}. Its recent article
\emph{Integer Pressure and a Missing Taylor Coefficient} proves
all-positive-integer phase analyticity, a vanishing $2m$th Taylor
coefficient, and a sixth-moment phase optimum \cite{IntegerRepo}.
That article explicitly excludes noninteger regularity. The present work
addresses this excluded direction rather than repeating its finite-matrix
construction or its optimization theorem.

The main observation is that the periodized sinc function appearing at
integer orders still makes sense at every real exponent $s>1$. Its
Fourier support need not be finite, but its transfer iterates can be
estimated explicitly. The loss of finite Fourier support is therefore
not the fundamental obstacle. The surviving fractional power in the
fixed-point eigenvalue equation is the actual source of nonanalyticity.

\subsection{Status and boundaries of the claim}
All central arguments below are supplied in ordinary mathematics. They
use no unverified theorem from the repository as an axiom. The repository
is credited for the integer-order predecessor, while the new proof
rederives the atomic eigenfunction and the coefficient cancellation in a
larger real-exponent framework.

The literature check included both primary papers above, the repository's
integer-pressure package, and targeted searches for fractional phase
singularities. No matching sharp classification was located in those
sources. This is a bounded search, not a guarantee of worldwide priority.
Nothing in this package is a Lean or Rocq proof, and numerical agreement
is not substituted for a proof of an infinite statement.

\begin{center}
\begin{tabular}{@{}p{.27\textwidth}p{.66\textwidth}@{}}
\toprule
\textbf{Status} & \textbf{Content}\\
\midrule
Prior work & Positive-order continuity; order-two analyticity; the repository's
integer-order analyticity and missing coefficient.\\[1.5mm]
Proved in this draft & The noninteger atomic cusp and exact integer differentiability
class; a sharp H\"older operator-norm relaxation law; the all-base extension;
a local uniform finite-size moment expansion.\\[1.5mm]
Not settled here & The critical and subcritical phase asymptotics, regularity at
arbitrary nonzero phases, endpoint H\"older regularity on whole neighborhoods,
and independent priority or formal verification.\\
\bottomrule
\end{tabular}
\end{center}

\section{Definitions and main results}
\label{sec:main}

\subsection{Digital masks and pressure}
Let $b\ge2$ be an integer and let $T_b(x)=bx\pmod1$ on
$\TT=\RR/\ZZ$. Define the normalized digital mask
\begin{equation}
 d_b(x)=\frac1b\sum_{r=0}^{b-1}e^{2\pi\ii rx},\qquad
 m_b(x)=|d_b(x)|
 =\left|\frac{\sin\pi bx}{b\sin\pi x}\right|,
 \label{eq:mask}
\end{equation}
with removable values at integers. For $s>1$ and a real phase $c$, put
\begin{equation}
 w_{b,s,c}(x)=m_b(x-c)^s,\qquad
 (\Lop_{b,s,c}f)(x)=\sum_{k=0}^{b-1}
 w_{b,s,c}\left(\frac{x+k}{b}\right)f\left(\frac{x+k}{b}\right).
 \label{eq:operator}
\end{equation}
There is no factor $1/b$ in this transfer operator.

Let $\mathcal D_{b,n}$ be the $b^n$ closed $b$-adic intervals of length
$b^{-n}$. Shared endpoints cause no problem for the integrals used below.
Set
\begin{align}
 F_{b,s,n,c}(x)&=\prod_{j=0}^{n-1}w_{b,s,c}(T_b^jx),\\
 Z_{b,s,n}(c)&=\sum_{I\in\mathcal D_{b,n}}\sup_{x\in I}F_{b,s,n,c}(x),\\
 P_{b,s}(c)&=\lim_{n\to\infty}\frac1n\log Z_{b,s,n}(c).
 \label{eq:pressure}
\end{align}
The existence and spectral interpretation of the last limit near $c=0$
will be proved, not assumed. All logarithms are natural.

For the binary potential in the introduction,
\begin{equation}
 p_q(c):=P(q\psi_c)=P_{2,2q}(c),\qquad
 \psi_c(x)=\log\cos^2\pi(x-c).
 \label{eq:binary-translation}
\end{equation}
Thus $s$ is the absolute-moment exponent, while $q=s/2$ is the usual
pressure order. This distinction accounts for the exponent $2q$ in the
binary theorem.

\subsection{The fractional normal form}
For $|c|<1/b$, define the positive analytic function
\begin{equation}
 D_b(c)=\frac{\sin\pi bc}{b\sin\pi c},\qquad D_b(0)=1,
 \qquad C_{b,s}=2(b^s-1)\zeta(s).
 \label{eq:constants}
\end{equation}

\begin{theorem}[Atomic-phase normal form]
\label{thm:normal}
Fix $b\ge2$ and $s>1$. There is $\varepsilon>0$ such that on
$|c|<\varepsilon$ the pressure exists, is even, and has the exact form
\begin{equation}
 P_{b,s}(c)=s\log D_b(c)+
 \log\bigl(1+|c|^s E_{b,s}(c)\bigr),
 \label{eq:normal}
\end{equation}
where $E_{b,s}$ is positive and even, satisfies
\begin{equation}
 E_{b,s}(0)=C_{b,s},
 \label{eq:Ezero}
\end{equation}
and is locally $C^r$ for every nonnegative integer $r<s$.
Consequently,
\begin{equation}
 P_{b,s}(c)=s\log D_b(c)+C_{b,s}|c|^s+o(|c|^{s+1}).
 \label{eq:asymp-strong}
\end{equation}
For even integer $s$, both the pressure and $E_{b,s}$ are real-analytic
near zero.
\end{theorem}

The one-extra-order remainder in \cref{eq:asymp-strong} is useful but is
not needed for the nonanalyticity conclusion. It follows from evenness
and one continuous derivative of $E_{b,s}$, not from an assumed power
series for that function.

\begin{theorem}[Sharp integer differentiability class]
\label{thm:regularity}
Let $s>1$ and assume that $s$ is not an even integer. Put
$k=\lceil s\rceil-1$. Then $P_{b,s}$ is $C^k$ in a neighborhood of zero,
but its derivative of order $k+1$ at zero does not exist. More precisely,
for every integer $0\le r<s$, as $c\to0$ through nonzero real values,
\begin{align}
 \frac{\mathrm d^r}{\mathrm dc^r}
 \bigl(P_{b,s}(c)-s\log D_b(c)\bigr)
 & = C_{b,s}(s)_r\,\sgn(c)^r |c|^{s-r}
   +o(|c|^{s-r}),
 \label{eq:derivative-cusp}\\
 (s)_r&=s(s-1)\cdots(s-r+1),\qquad (s)_0=1.
\end{align}
\end{theorem}

\begin{corollary}[Binary pressure]
\label{cor:binary}
For every noninteger $q>1/2$, $p_q$ is locally
$C^{\lceil2q\rceil-1}$ but has no derivative of order $\lceil2q\rceil$
at zero. Its expansion is
\begin{equation}
 p_q(c)=2q\log\cos\pi c+
 2(2^{2q}-1)\zeta(2q)|c|^{2q}+o(|c|^{2q+1}).
 \label{eq:binary-cusp}
\end{equation}
In particular, noninteger positive orders above $1/2$ do not admit a
real-analytic phase dependence on the whole circle.
\end{corollary}

\subsection{The sharp relaxation theorem}
Use the circle distance $\delta(x,y)$ and the norm
\[
 \|f\|_{\alpha}=\|f\|_\infty+[f]_{\alpha},\qquad
 [f]_{\alpha}=\sup_{x\ne y}\frac{|f(x)-f(y)|}{\delta(x,y)^\alpha}
\]
on the complex Banach space $C^\alpha(\TT)$, where $0<\alpha\le1$.
The case $\alpha=1$ means Lipschitz functions, not continuously
differentiable functions. Define
\begin{equation}
 H_s(x)=\sum_{\ell\in\ZZ}|\sinc(x+\ell)|^s,
 \qquad \sinc u=\frac{\sin\pi u}{\pi u},\quad \sinc0=1,
 \qquad \Proj_sf=f(0)H_s.
 \label{eq:H}
\end{equation}

\begin{theorem}[Sharp atomic relaxation]
\label{thm:relax}
For every $b\ge2$, $s>1$, and $0<\alpha\le1$, $H_s$ is strictly
positive and belongs to $C^1(\TT)$, and
\begin{equation}
 \Lop_{b,s,0}H_s=H_s,\qquad
 \Lop_{b,s,0}\Proj_s=\Proj_s\Lop_{b,s,0}=\Proj_s.
\end{equation}
There are positive constants $a,A$, depending only on $b,s,\alpha$, such
that for all $n\ge1$,
\begin{equation}
 aR_{b,s,\alpha}(n)\le
 \|\Lop_{b,s,0}^{\,n}-\Proj_s\|_{C^\alpha\to C^\alpha}
 \le AR_{b,s,\alpha}(n),
 \label{eq:sharp-relax}
\end{equation}
where
\begin{equation}
 R_{b,s,\alpha}(n)=
 \begin{cases}
 b^{-(s-1)n},&1<s<1+\alpha,\\
 (n+1)b^{-\alpha n},&s=1+\alpha,\\
 b^{-\alpha n},&s>1+\alpha.
 \end{cases}
 \label{eq:rate}
\end{equation}
In particular, $1$ is an algebraically simple isolated eigenvalue and all
remaining spectrum lies in
$|z|\le\max\{b^{-\alpha},b^{1-s}\}<1$.
\end{theorem}

The logarithmic factor at $s=1+\alpha$ is necessary in this operator norm.
It is detected by the single test function $\delta(x,0)^\alpha$; it is
not an artifact of an abstract spectral estimate.

\section{The atomic eigenfunction and the telescoping identity}
\label{sec:atomic}

\subsection{Construction and regularity of the periodization}
For $x\notin\ZZ$, the series in \cref{eq:H} becomes
\begin{equation}
 H_s(x)=\frac{|\sin\pi x|^s}{\pi^s}
 \sum_{\ell\in\ZZ}|x+\ell|^{-s}.
 \label{eq:Hsum}
\end{equation}
On compact subsets away from the integers, the series and its first
derivative converge uniformly. Near zero, separate the $\ell=0$ term:
\begin{equation}
 H_s(x)=|\sinc x|^s+
 \frac{|\sin\pi x|^s}{\pi^s}
 \sum_{\ell\ne0}|x+\ell|^{-s}.
 \label{eq:Hsplit}
\end{equation}
The second sum is continuously differentiable there, with uniformly
convergent derivative series. Since $s>1$, multiplication by
$|\sin\pi x|^s$ preserves $C^1$ regularity at zero. Thus $H_s\in C^1$,
$H_s(0)=1$, and periodicity follows by reindexing.

Every $x$ has an integer translate $u$ with $|u|\le1/2$. The corresponding
summand is at least $(2/\pi)^s$. Hence
\begin{equation}
 H_s(x)\ge(2/\pi)^s>0\qquad(x\in\TT).
 \label{eq:Hpositive}
\end{equation}
This uniform lower bound is important: it allows a small perturbation of
the eigenfunction to remain strictly positive.

\subsection{The eigenfunction identity}
For $y=(x+k)/b$ with $0\le k<b$, use
$|\sin\pi by|=|\sin\pi x|$ in \cref{eq:Hsum}. Away from removable
points,
\begin{align*}
 m_b(y)^sH_s(y)
 &=\frac{|\sin\pi x|^s}{b^s\pi^s}
   \sum_{\ell\in\ZZ}|y+\ell|^{-s}\\
 &=\frac{|\sin\pi x|^s}{\pi^s}
   \sum_{\ell\in\ZZ}|x+k+b\ell|^{-s}.
\end{align*}
Summing over $k$ partitions $\ZZ$ into its residue classes modulo $b$,
giving $\Lop_{b,s,0}H_s=H_s$. Continuity handles the removable points.
At $x=0$, only the inverse branch $y=0$ has nonzero weight, so
\begin{equation}
 (\Lop_{b,s,0}f)(0)=f(0).
 \label{eq:delta-left}
\end{equation}
Together these identities prove the projector relations in
\cref{thm:relax}.

\subsection{An exact formula for every iterate}
Let $N=b^n$. The product of the masks telescopes:
\begin{equation}
 \prod_{j=0}^{n-1}m_b(b^jy)^s
 =\left|\frac{\sin\pi Ny}{N\sin\pi y}\right|^s.
 \label{eq:telescope}
\end{equation}
Therefore, for any complete set $K_N$ of $N$ consecutive integer residues
modulo $N$,
\begin{equation}
 (\Lop_{b,s,0}^{\,n}f)(x)
 =\sum_{k\in K_N}a_{N,k}(x)f\left(\frac{x+k}{N}\right),\qquad
 a_{N,k}(x)=\left|
 \frac{\sin\pi x}{N\sin(\pi(x+k)/N)}\right|^s.
 \label{eq:iterate}
\end{equation}
Take representatives $x\in[-1/2,1/2]$ and a centered $K_N$ containing
zero, for example $K_N=\{-N/2,\ldots,N/2-1\}$ when $N$ is even.
At the removable zero of $a_{N,0}$, set its value to $1$.

This formula is the analytic substitute for a finite moment matrix.
The proof below keeps track of both the total mass and its H\"older
variation, which is essential for a genuine spectral perturbation argument.

\section{Proof of the sharp relaxation law}
\label{sec:relax-proof}

\subsection{Uniform bounds on the branch coefficients}
\begin{lemma}
\label{lem:coeff}
For fixed $s>1$, there is a constant $C_s$ independent of $N\ge2$ such
that on $[-1/2,1/2]$,
\begin{equation}
 \|a_{N,0}\|_{C^1}\le C_s,\qquad
 \|a_{N,k}\|_{C^1}\le C_s|k|^{-s}\quad(0\ne k\in K_N).
 \label{eq:coeff-bound}
\end{equation}
Here $C^1$ is used only for the scalar coefficient functions on the
interval.
\end{lemma}
\begin{proof}
For $k=0$,
\[
 a_{N,0}(x)=\left|\frac{\sinc x}{\sinc(x/N)}\right|^s.
\]
Both factors are uniformly bounded away from zero on the relevant
intervals, and their first derivatives are uniformly bounded.

For $k\ne0$, one has $|x+k|\ge|k|/2$. The quantity
$(x+k)/N$ stays in a fixed interval strictly inside $(-1,1)$, and thus
\[
 N|\sin(\pi(x+k)/N)|\ge c_0|x+k|\ge(c_0/2)|k|
\]
with an absolute positive $c_0$. Differentiating the quotient, the derivative
of the numerator power is uniformly bounded because $s>1$. The logarithmic
derivative of the denominator is bounded by $C/|x+k|$:
\[
 \left|\frac{\pi}{N}\cot\frac{\pi(x+k)}N\right|
 \le\frac{C}{|x+k|}.
\]
These estimates give \cref{eq:coeff-bound}.
\end{proof}

\subsection{Upper bound on the invariant ideal}
Let
\[
 \mathcal I_\alpha=\{f\in C^\alpha(\TT):f(0)=0\}.
\]
By \cref{eq:delta-left}, this is an invariant closed subspace. For
$f\in\mathcal I_\alpha$,
\begin{equation}
 |f(y)|\le[f]_\alpha\delta(y,0)^\alpha.
 \label{eq:vanishing-bound}
\end{equation}
Put $f_{N,k}(x)=f((x+k)/N)$. On the representative interval,
\begin{align}
 [f_{N,k}]_\alpha&\le N^{-\alpha}[f]_\alpha,\\
 \|f_{N,0}\|_\infty&\le C N^{-\alpha}[f]_\alpha,\\
 \|f_{N,k}\|_\infty&\le C N^{-\alpha}|k|^\alpha[f]_\alpha
 \quad(k\ne0).
 \label{eq:branch-f}
\end{align}
The estimates remain valid when an inverse image crosses a chosen
representative endpoint, because $f$ is periodic and circle distance is
at most the distance between its lifts.

Using the product inequality
$[uv]_\alpha\le\|u\|_\infty[v]_\alpha+[u]_\alpha\|v\|_\infty$,
\cref{lem:coeff}, and \cref{eq:branch-f}, we obtain
\begin{equation}
 \|\Lop_{b,s,0}^{\,n}f\|_\alpha
 \le C N^{-\alpha}
 \left(1+\sum_{k=1}^{N}k^{\alpha-s}\right)\|f\|_\alpha.
 \label{eq:sum-bound}
\end{equation}
For completeness, the interval estimate gives the circle estimate up to a
fixed factor: for two points separated by the cut, split their difference
at the identified endpoints and use
$a^\alpha+b^\alpha\le2^{1-\alpha}(a+b)^\alpha$.

The elementary power-sum alternatives are
\begin{equation}
 1+\sum_{k=1}^{N}k^{\alpha-s}\asymp
 \begin{cases}
 N^{1+\alpha-s},&s<1+\alpha,\\
 1+\log N,&s=1+\alpha,\\
 1,&s>1+\alpha.
 \end{cases}
 \label{eq:powersum}
\end{equation}
Since $N=b^n$, \cref{eq:sum-bound} gives the upper rate in
\cref{eq:sharp-relax} on $\mathcal I_\alpha$.
For arbitrary $f$, write
\[
 f=f(0)H_s+g,\qquad g\in\mathcal I_\alpha,
 \qquad \|g\|_\alpha\le C_s\|f\|_\alpha.
\]
Because $H_s$ is fixed by the operator, the same bound holds for
$\Lop^n-\Proj_s$ on the whole space.

\subsection{Matching lower bound}
\begin{lemma}
\label{lem:lower}
The fixed function $f_\alpha(x)=\delta(x,0)^\alpha$ satisfies
\begin{equation}
 \left(\Lop_{b,s,0}^{\,n}f_\alpha\right)(1/2)
 \ge c N^{-\alpha}
 \left(1+\sum_{k=1}^{\lfloor N/4\rfloor}k^{\alpha-s}\right)
 \label{eq:lower}
\end{equation}
for all sufficiently large $N=b^n$.
\end{lemma}
\begin{proof}
At $x=1/2$, the numerator of $a_{N,k}$ has absolute value $1$.
For $0\le k\le N/4$, the inverse image $(k+1/2)/N$ lies in $[0,1/2]$,
so its distance to zero is exactly $(k+1/2)/N$. Also,
\[
 N\sin\frac{\pi(k+1/2)}N\le\pi(k+1/2).
\]
Consequently each such summand in \cref{eq:iterate} is at least
\[
 \pi^{-s}N^{-\alpha}(k+1/2)^{\alpha-s}.
\]
All summands are nonnegative, and the $k=0$ term supplies the initial
constant in \cref{eq:lower}.
\end{proof}

The function $f_\alpha$ belongs to $C^\alpha$ with bounded norm and
$f_\alpha(0)=0$. Thus $\Proj_sf_\alpha=0$. Applying
\cref{eq:powersum} to \cref{eq:lower} gives the lower operator-norm bound.
Changing constants covers the finitely many small $n$.
This proves the two-sided estimate in \cref{thm:relax}.

\subsection{The spectral consequence}
Write $R=\Lop_{b,s,0}-\Proj_s$. The projector relations imply
$R^n=\Lop_{b,s,0}^{\,n}-\Proj_s$ for $n\ge1$, and the estimates give
\[
 r(R)\le\max\{b^{-\alpha},b^{1-s}\}<1.
\]
Since the range of $\Proj_s$ is one-dimensional, this decomposes the
operator into the simple eigenvalue $1$ and a strictly smaller spectral
part. In particular there is no generalized eigenvector attached to $1$.
This supplies the spectral gap needed below directly, rather than by
invoking a positivity theorem that might fail at an atomic phase.

\section{Parameter regularity and spectral perturbation}
\label{sec:perturb}

\subsection{A translation lemma on H\"older spaces}
\begin{lemma}
\label{lem:translation}
Fix $s>1$ and a nonnegative integer $r<s$. Choose
$0<\alpha<\min\{1,s-r\}$. Then
\[
 c\longmapsto w_{b,s,c}\in C^\alpha(\TT),\qquad
 c\longmapsto\Lop_{b,s,c}\in\mathcal B(C^\alpha(\TT))
\]
are $C^r$ maps. If $s$ is an even integer, the maps are real-analytic.
\end{lemma}
\begin{proof}
The nontrivial zeros of the mask occur at $j/b$, $1\le j<b$, and are
simple zeros before taking the absolute value. In a local coordinate $u$
around each zero, the weight has the form $|u|^s v(u)$ with $v$ positive
and real-analytic. Away from those zeros it is real-analytic.
Its $r$th derivative is H\"older continuous of every exponent
$\gamma<\min\{1,s-r\}$. Choose such a $\gamma$ larger than $\alpha$.

Translations of a $C^\gamma$ function are continuous in the $C^\alpha$
norm. One elementary proof splits the H\"older quotient into distances
at most $\eta$ and greater than $\eta$: the first part is bounded by
$C\eta^{\gamma-\alpha}$ and the second by
$2\eta^{-\alpha}$ times the uniform translation error. First choose
$\eta$, then let the translation tend to zero.
Applying this to the $r$th derivative, and using the integral form of the
Taylor remainder for the lower derivatives, proves $C^r$ translation
dependence in $C^\alpha$.

Multiplication is bounded on $C^\alpha$, and the finite sum over inverse
branches is bounded there. Hence the same parameter regularity holds for
the operator. When $s=2m$, the weight is a trigonometric polynomial, so
its translation depends analytically on $c$ in the operator norm.
\end{proof}

The strict inequality $\alpha<s-r$ is intentional. Translating a function
at its exact H\"older endpoint need not be norm-continuous. Passing to a
slightly weaker H\"older space avoids that pitfall without losing an
integer derivative of the pressure.

\subsection{The isolated eigenvalue persists}
\begin{proposition}
\label{prop:perron}
For every fixed $b\ge2$, $s>1$, and integer $r<s$, a neighborhood of
zero has $C^r$ functions $\rho(c)>0$ and $h_c\in C^\alpha(\TT)$ such that
\begin{equation}
 \Lop_{b,s,c}h_c=\rho(c)h_c,\qquad
 h_c(0)=1,\qquad h_c>0,\qquad
 \rho(0)=1,\quad h_0=H_s.
 \label{eq:perron}
\end{equation}
The eigenvalue is algebraically simple and strictly dominant. There is
also a $C^r$ rank-one spectral projector $\Proj_c$, and, after possibly
shrinking the neighborhood,
\begin{equation}
 \|\rho(c)^{-n}\Lop_{b,s,c}^{\,n}-\Proj_c\|_{\alpha\to\alpha}
 \le K\theta^n\qquad(n\ge0)
 \label{eq:uniform-gap}
\end{equation}
with $0<\theta<1$ and $K<\infty$ independent of $c$. For even integer
$s$, the spectral data are real-analytic.
\end{proposition}
\begin{proof}
By \cref{thm:relax}, a small circle $\Gamma$ around $1$ separates that
simple eigenvalue from the rest of the spectrum of $\Lop_{b,s,0}$.
Operator-norm continuity and the Neumann series preserve the resolvent
on $\Gamma$ for small $c$. Define
\[
 \Proj_c=\frac1{2\pi\ii}\int_\Gamma
 (z-\Lop_{b,s,c})^{-1}\dd z.
\]
Differentiating the resolvent identity gives the stated $C^r$ dependence.
The projectors remain rank one: sufficiently close projections have
isomorphic ranges, by restricting one projector to the other's range.
Since $\Proj_0 1=H_s$, set
\[
 h_c=\frac{\Proj_c1}{(\Proj_c1)(0)}.
\]
The denominator tends to $1$, and \cref{eq:Hpositive} ensures that
$h_c$ remains strictly positive. The invariant one-dimensional range
supplies an eigenvalue; evaluating at zero gives
$\rho(c)=(\Lop_{b,s,c}h_c)(0)$. Reality, positivity, and $C^r$ regularity
follow immediately.

The remaining spectrum stays in a disk of radius less than $1$, whereas
$\rho(c)$ tends to $1$. This can be seen directly from resolvent stability:
outside fixed neighborhoods of the old spectrum, cover a bounded region
by finitely many resolvent neighborhoods and use the Neumann series;
large $|z|$ are handled by the uniform operator-norm bound.
Thus the eigenvalue remains strictly dominant.

For the uniform estimate, normalize the complementary operator by
$\rho(c)$. A fixed circle $|z|=\theta<1$ encloses all of its spectra for
small $c$, with uniformly bounded resolvents on that circle. The Cauchy
integral formula for its powers gives \cref{eq:uniform-gap}. Enlarge $K$
to include $n=0$. Analytic input gives analytic resolvents and spectral
data when $s$ is even.
\end{proof}

Although different $r$ may require different H\"older spaces, they give
the same positive eigenvalue and normalized positive eigenfunction near
zero. To see uniqueness without comparing the Banach spaces, compare
two strictly positive eigenfunctions by their maximum and minimum ratio.
Positivity of $\Lop$ forces their positive eigenvalues to be equal; their
difference then lies in the unique eigenspace in a common weaker
H\"older space. Hence the resulting pressure regularities are compatible.

\subsection{Reflection symmetry}
Let $(Jf)(x)=f(-x)$. Evenness of the mask gives
\[
 J\Lop_{b,s,c}=\Lop_{b,s,-c}J.
\]
Uniqueness and the normalization at zero imply
\begin{equation}
 \rho(-c)=\rho(c),\qquad h_{-c}(x)=h_c(-x).
 \label{eq:reflection}
\end{equation}
The atomic eigenfunction is therefore the common limit of the perturbed
positive eigenfunctions from both sides.

\section{Why the spectral eigenvalue is the pressure}
\label{sec:pressure}

A positive eigenvalue on a chosen function space does not by itself
identify cylinder-supremum pressure. The following comparison closes
that gap.

\begin{lemma}[A polynomial variation bound]
\label{lem:bernstein}
For every $p\ge1$ there is a constant $B_p$ such that every trigonometric
polynomial $Q$ of degree at most $D\ge1$ satisfies
\begin{equation}
 \int_0^1\left|\frac{\mathrm d}{\mathrm dx}|Q(x)|^p\right|\dd x
 \le B_pD\int_0^1|Q(x)|^p\dd x.
 \label{eq:variation}
\end{equation}
The derivative is interpreted almost everywhere when necessary.
\end{lemma}
\begin{proof}
A convolution proof of the elementary estimate
$\|Q'\|_p\le CD\|Q\|_p$ is given in \cref{app:bernstein}.
For $p>1$, the chain rule and H\"older's inequality give
\[
 \int\bigl|(|Q|^p)'\bigr|
 \le p\int|Q|^{p-1}|Q'|
 \le p\|Q\|_p^{p-1}\|Q'\|_p.
\]
For $p=1$, absolute continuity and $\bigl|(|Q|)'\bigr|\le|Q'|$
almost everywhere give the same conclusion.
\end{proof}

Define the digital product
\begin{equation}
 Q_{b,n,c}(x)=\prod_{j=0}^{n-1}
 \left(\sum_{r=0}^{b-1}e^{2\pi\ii r(b^jx-c)}\right).
 \label{eq:digital-product}
\end{equation}
Its degree in $e^{2\pi\ii x}$ is $b^n-1$, and
\begin{equation}
 F_{b,s,n,c}(x)=b^{-sn}|Q_{b,n,c}(x)|^s.
 \label{eq:Qnormal}
\end{equation}
On an interval of length $1/N$, the supremum of a nonnegative absolutely
continuous function is at least its average and at most its average plus
its variation. Summing over the partition, with $N=b^n$, gives
\begin{equation}
 N\int F\le Z_{b,s,n}(c)
 \le N\int F+\int|F'|
 \le(1+B_s)N\int F.
 \label{eq:Zsandwich}
\end{equation}
The constants do not depend on $n$ or the phase.

Changing variables in the inverse-branch sums yields
\begin{equation}
 \int_0^1\Lop_{b,s,c}^{\,n}1\dd x=N\int_0^1F_{b,s,n,c}(x)\dd x.
 \label{eq:int-transfer}
\end{equation}
By \cref{prop:perron}, the left side equals
\[
 \rho(c)^n\bigl(\mathcal A_{b,s}(c)+O(\theta^n)\bigr),\qquad
 \mathcal A_{b,s}(c)=\int_0^1\Proj_c1\dd x.
\]
At zero, $\mathcal A_{b,s}(0)=\int H_s>0$. Continuity gives a uniform
positive lower bound on a smaller neighborhood. Taking logarithms in
\cref{eq:Zsandwich} proves both existence of the limit and
\begin{equation}
 P_{b,s}(c)=\log\rho(c).
 \label{eq:pressure-rho}
\end{equation}
No bounded-distortion estimate for the logarithm of a vanishing weight
has been used.

\section{The fixed-point identity and the fractional cusp}
\label{sec:cusp}

\subsection{An exact scalar equation}
Evaluate the eigenfunction equation at $x=0$. The principal inverse
branch gives $D_b(c)^s$, since $h_c(0)=1$. Every other inverse branch
has numerator $|\sin\pi bc|^s$. Hence
\begin{equation}
 \rho(c)=D_b(c)^s+
 \frac{|\sin\pi bc|^s}{b^s}
 \sum_{k=1}^{b-1}
 \frac{h_c(k/b)}{|\sin\pi(k/b-c)|^s}.
 \label{eq:scalar}
\end{equation}
Choose the phase neighborhood small enough that the denominators do not
vanish. Write
\begin{equation}
 B_{b,s}(c)=\sum_{k=1}^{b-1}
 \frac{h_c(k/b)}{|\sin\pi(k/b-c)|^s}.
 \label{eq:B}
\end{equation}
Dividing \cref{eq:scalar} by $D_b(c)^s$ gives the particularly simple
identity
\begin{equation}
 \rho(c)=D_b(c)^s\bigl(1+|\sin\pi c|^sB_{b,s}(c)\bigr).
 \label{eq:factor}
\end{equation}
Thus the logarithmic singularity of the original potential has become
an explicit fractional power in a scalar eigenvalue equation.

The function $B_{b,s}$ is positive and $C^r$ for every integer $r<s$.
By \cref{eq:reflection}, replacing $k$ by $b-k$ shows that it is even.
Define
\begin{equation}
 E_{b,s}(c)=\left(\frac{\sin\pi c}{c}\right)^sB_{b,s}(c),
 \qquad E_{b,s}(0)=\pi^sB_{b,s}(0),
 \label{eq:E}
\end{equation}
where the quotient is positive in the small neighborhood under
consideration. Its power is real-analytic there. Equations
\cref{eq:pressure-rho,eq:factor} now give the exact normal form.

\subsection{Evaluation of the cusp coefficient}
Using \cref{eq:Hsum},
\begin{align}
 \pi^sB_{b,s}(0)
 &=\sum_{k=1}^{b-1}\sum_{\ell\in\ZZ}|k/b+\ell|^{-s}\\
 &=b^s\sum_{\substack{j\in\ZZ\\b\nmid j}}|j|^{-s}
 =2(b^s-1)\zeta(s).
 \label{eq:amplitude}
\end{align}
All rearrangements are justified by absolute convergence, since $s>1$.
The same coefficient works in every base, with the base entering only
through the omitted multiples in the final zeta sum.

Since $E_{b,s}$ is even and $C^1$, its derivative at zero is zero, so
$E_{b,s}(c)=C_{b,s}+o(|c|)$. Consequently,
\begin{align*}
 \log(1+|c|^sE_{b,s}(c))
 &=C_{b,s}|c|^s+o(|c|^{s+1})+O(|c|^{2s})\\
 &=C_{b,s}|c|^s+o(|c|^{s+1}),
\end{align*}
because $s>1$. This proves \cref{thm:normal}, including its stronger
remainder.

\subsection{Differentiation and the exact loss of smoothness}
Fix an integer $r<s$. On either side of zero, differentiate
\cref{eq:normal}. The leading product term is
\[
 C_{b,s}\frac{\mathrm d^r}{\mathrm dc^r}|c|^s
 =C_{b,s}(s)_r\sgn(c)^r|c|^{s-r}.
\]
Every term with a positive number $j$ of derivatives falling on $E$
has an additional factor $|c|^j$ relative to this scale and a bounded
derivative $E^{(j)}$. The term with no derivative falling on $E$ has
$E(c)-C_{b,s}=o(1)$. Finally, the nonlinear part of $\log(1+u)-u$
contributes $O(|c|^{2s-r})$ under these differentiations. This proves
\cref{eq:derivative-cusp}.

Let $k=\lceil s\rceil-1$ and subtract the analytic function
$a(c)=s\log D_b(c)$. The derivatives of $P-a$ up to order $k$ exist
continuously and vanish at zero. If $s$ is not an integer, then
$0<s-k<1$, so
\[
 \frac{(P-a)^{(k)}(c)-(P-a)^{(k)}(0)}{c}
 =C_{b,s}(s)_k\sgn(c)^{k-1}|c|^{s-k-1}
   +o(|c|^{s-k-1})
\]
has no finite limit. If $s$ is an odd integer, $k=s-1$ is even and the
two one-sided limits are $\pm C_{b,s}s!$. Thus a derivative of order
$k+1$ fails in both cases. The only remaining integer case is even $s$,
which is analytic and is treated next. This proves
\cref{thm:regularity,cor:binary}.

\section{Integer resonance and explicit consequences}
\label{sec:integer}

\subsection{Why the integer coefficient disappears}
Euler's sine product, followed by an absolutely convergent logarithmic
expansion, gives for $|c|<1/b$
\begin{equation}
 \log D_b(c)
 =-\sum_{j\ge1}\frac{(b^{2j}-1)\zeta(2j)}{j}c^{2j}.
 \label{eq:log-mask}
\end{equation}
For $s=2m$, \cref{eq:normal} is analytic. At degree $2m$, the first
term contributes
\[
 -\frac{2m}{m}(b^{2m}-1)\zeta(2m)
 =-2(b^{2m}-1)\zeta(2m),
\]
and the second contributes exactly the opposite number, by
\cref{eq:amplitude}. Hence:

\begin{corollary}[All-base integer cancellation]
\label{cor:cancel}
For every integer $m\ge1$ and base $b\ge2$,
\begin{equation}
 [c^{2m}]P_{b,2m}(c)=0,
 \label{eq:cancel}
\end{equation}
and for $1\le j<m$,
\begin{equation}
 [c^{2j}]P_{b,2m}(c)
 =-\frac{2m}{j}(b^{2j}-1)\zeta(2j).
 \label{eq:lowercoeff}
\end{equation}
\end{corollary}

The binary cancellation is credited to the repository predecessor
\cite{IntegerRepo}; the role here is to explain it as the even-integer
specialization of a real-exponent normal form and extend it to arbitrary
integer bases. At a noninteger exponent, there is no analytic coefficient
of degree $s$ available to cancel $|c|^s$.

\subsection{The order-one normalization}
When $s=2$, the pressure vanishes identically, not merely to a finite
Taylor order. Indeed, discrete Fourier orthogonality gives
\begin{equation}
 \sum_{k=0}^{b-1}\left|
 \frac1b\sum_{r=0}^{b-1}e^{2\pi\ii r((x+k)/b-c)}\right|^2=1.
 \label{eq:g-normalization}
\end{equation}
To check it, expand the square and sum over $k$. Only $r=t$ survives,
since $|r-t|<b$. Thus $\Lop_{b,2,c}1=1$ for every phase. The pressure
comparison in \cref{sec:pressure} gives $P_{b,2}(c)=0$ everywhere.
This is consistent with the normal form and isolates the simplest
integer resonance.

\subsection{Local sign and curvature}
For $1<s<2$, the fractional term dominates the analytic quadratic term,
so
\begin{equation}
 P_{b,s}(c)\sim C_{b,s}|c|^s>0\qquad(c\to0,\ c\ne0).
 \label{eq:sub-two}
\end{equation}
For $s>2$, the quadratic term dominates:
\begin{equation}
 P_{b,s}(c)
 =-\frac{s(b^2-1)\pi^2}{6}c^2+o(c^2),\qquad
 P_{b,s}''(0)=-\frac{s(b^2-1)\pi^2}{3}.
 \label{eq:curvature}
\end{equation}
Thus the atomic phase is a strict local minimum for $1<s<2$, a flat
normalization point for $s=2$, and a strict local maximum for $s>2$.
These assertions are local; no global optimization theorem is inferred.

For the binary family, several representative cases are:
\begin{center}
\begin{tabular}{@{}lll@{}}
\toprule
$q$ & Leading fractional coefficient & Exact local integer class\\
\midrule
$3/4$ & $2(2^{3/2}-1)\zeta(3/2)$ & $C^1$, no second derivative at $0$\\
$5/4$ & $2(2^{5/2}-1)\zeta(5/2)$ & $C^2$, no third derivative at $0$\\
$3/2$ & $14\zeta(3)$ & $C^2$, no third derivative at $0$\\
$7/4$ & $2(2^{7/2}-1)\zeta(7/2)$ & $C^3$, no fourth derivative at $0$\\
$2$ & The degree-four term cancels & Real-analytic near $0$\\
\bottomrule
\end{tabular}
\end{center}
For example,
\[
 p_{3/2}(c)=-\frac32\pi^2c^2+14\zeta(3)|c|^3+o(|c|^3).
\]
The odd absolute power is a genuine two-sided cusp even though its
exponent is an integer.

\section{Finite-size moments and the critical boundary}
\label{sec:moments}

\subsection{Uniform asymptotics for finite digital products}
\begin{corollary}
\label{cor:moments}
For fixed $b\ge2$ and $s>1$, there is a phase neighborhood of zero on
which
\begin{equation}
 \int_0^1|Q_{b,n,c}(x)|^s\dd x
 =\mathcal A_{b,s}(c)
 \exp\!\left(n\bigl((s-1)\log b+P_{b,s}(c)\bigr)\right)
 \bigl(1+O(\theta^n)\bigr),
 \label{eq:moment-asymp}
\end{equation}
uniformly in the phase, where $0<\theta<1$ and
$\mathcal A_{b,s}(c)>0$. The amplitude is $C^r$ for each integer $r<s$,
and
\begin{equation}
 \mathcal A_{b,s}(0)=\int_0^1H_s(x)\dd x
 =\int_{\RR}|\sinc u|^s\dd u.
 \label{eq:amplitude-moment}
\end{equation}
\end{corollary}
\begin{proof}
Combine \cref{eq:uniform-gap,eq:int-transfer,eq:Qnormal}. The leading
coefficient stays bounded away from zero, so an absolute error may be
written as a relative error. The final equality follows from Tonelli's
theorem applied to the nonnegative periodization; the integral converges
because $s>1$.
\end{proof}

At $c=0$, the finite product is simply
$\sum_{r=0}^{b^n-1}e^{2\pi\ii rx}$. Away from zero, it is a
phase-deformed digit product; for $b=2$, $c=1/2$, its coefficients are
the classical Thue--Morse signs. The theorem is local near the atomic
phase and does not supply a uniform estimate reaching $c=1/2$.

For binary Riesz-product measures, the established pressure--$L^q$
spectrum identity identifies the two quantities up to division by
$\log2$ at nonzero phases \cite[Theorem~2.6]{GKS}. At orders $q>1/2$,
the pressure tends to zero at the atomic phase, consistent with the
trivial spectrum of the point mass. Consequently the same local cusp
classification transfers to that phase-dependent spectrum, with its
coefficient divided by $\log2$. This corollary uses the cited
measure-theoretic theorem; the pressure proof above does not.

\subsection{Why the threshold is not a technical convenience}
The coefficient $C_{b,s}$ diverges as $s\downarrow1$, and the slow
relaxation factor $b^{1-s}$ tends to $1$. Both facts warn against passing
to the critical exponent inside the supercritical expansion.
There is also a direct critical obstruction.

\begin{proposition}[Critical growth at the atomic phase]
\label{prop:critical}
Extend the definition of the continuous weight to $s=1$. Then
\begin{equation}
 \|\Lop_{b,1,0}^{\,n}1\|_\infty\asymp n+1.
 \label{eq:critical-growth}
\end{equation}
In particular, the rank-one bounded-power convergence used above cannot
hold at the critical exponent.
\end{proposition}
\begin{proof}
The scalar iterate formula \cref{eq:iterate} is still valid. Its $k=0$
coefficient is uniformly bounded, and the remaining coefficients are
bounded by $C/|k|$. Hence the supremum is at most $C(1+\log N)$.
At $x=1/2$, the terms with $1\le k\le N/4$ are at least
$1/(\pi(k+1/2))$, giving a matching lower bound. Use $N=b^n$.
\end{proof}

This proposition does not determine the phase asymptotic of $P_{b,1}(c)$.
It identifies the failure of the precise mechanism used for $s>1$ and
motivates a separate critical analysis rather than an unjustified
continuation of the zeta coefficient.

\section{Computation, reproducibility, and verification limits}
\label{sec:checks}

\subsection{What the supplied program checks}
The executable \code{code/verify.py} carries out three independent kinds
of finite checks. It evaluates the atomic eigenfunction identity at
65-decimal-digit precision using
\begin{equation}
 H_s(x)=\left(\frac{\sin\pi x}{\pi}\right)^s
 \bigl(\zeta(s,x)+\zeta(s,1-x)\bigr),\qquad 0<x<1,
 \label{eq:hurwitz}
\end{equation}
it verifies the residue-class formula for $C_{b,s}$, and it checks the
integer cancellation after factoring out the common zeta value using
exact symbolic arithmetic. It also iterates a positive piecewise-linear
collocation of the transfer operator to test the cusp numerically.

The recorded run contains 112 high-precision eigenfunction checks and
28 high-precision amplitude checks for bases $2,3,5,7$ and seven real
exponents. Their largest observed absolute residuals are approximately
$1.2\times10^{-64}$ and $3.9\times10^{-61}$, respectively. The symbolic
cancellation check covers $m=1,\ldots,20$ with symbolic base $b$.
These computations test formulas; the all-exponent and all-base results
are supplied by the proofs, not inferred from those finite ranges.

\subsection{A stable numerical observable}
Directly subtracting $s\log D_b(c)$ from a numerical pressure loses
precision at large $s$ and small $c$. Instead the program evaluates the
normal-form excess with \code{log1p}:
\begin{equation}
 \mathcal R_{b,s}(c)=
 \frac{P_{b,s}(c)-s\log D_b(c)}{|c|^s}
 =\frac{\log(1+|\sin\pi c|^sB_{b,s}(c))}{|c|^s}.
 \label{eq:numeric-ratio}
\end{equation}
The theorem predicts $\mathcal R_{b,s}(c)\to C_{b,s}$. The following
values use the default grid, rounded for display.

\begin{center}
\begin{tabular}{@{}rrrrrr@{}}
\toprule
$b$ & $s$ & $c=10^{-2}$ & $c=3\cdot10^{-3}$ & $c=10^{-3}$ & $C_{b,s}$\\
\midrule
2 & 1.5 & 8.909155 & 9.418009 & 9.522764 & 9.553076\\
2 & 2.5 & 12.609818 & 12.505178 & 12.495476 & 12.494221\\
2 & 3.0 & 17.065680 & 16.850219 & 16.831182 & 16.828797\\
2 & 3.5 & 23.687926 & 23.281687 & 23.246062 & 23.241609\\
3 & 1.5 & 19.960605 & 21.496843 & 21.827067 & 21.923850\\
3 & 2.5 & 39.702426 & 39.194272 & 39.146656 & 39.140459\\
\bottomrule
\end{tabular}
\end{center}

The program records 18 such phase tests and two doubled-grid checks.
The grid has 65,536 points for base two and is rounded upward to a
multiple of the base when necessary. It uses nonnegative linear
interpolation, preserving positivity of the discretized operator.
The iteration is normalized by its value at zero, matching the proof's
eigenfunction normalization.

\subsection{What these numbers do not certify}
The floating-point residual is the residual of the discretized iteration,
not an enclosure of the true operator spectrum. No directed-rounding
interval certificate or rigorous discretization-error bound is claimed.
The doubled-grid tests are stability diagnostics only. The high-precision
zeta checks are also numerical rather than exact proofs; only the
separate symbolic cancellation identities use exact arithmetic.

The package deliberately separates these layers in
\code{results/verification.json}. It also includes a source ledger,
rebuild instructions, and a claims ledger. No repository file has been
modified by preparing this package.

\section{Further research questions}
\label{sec:questions}

\begin{question}[Critical phase asymptotics]
Determine the leading behavior of $P_{b,1}(c)$ as $c\to0$. The linear
critical growth in \cref{prop:critical}, the pole of $\zeta(s)$, and the
closing spectral gap should be treated together. An induced or renewal
operator may replace the lost rank-one limit. A guessed logarithmic law
is not asserted here.
\end{question}

\begin{question}[The subcritical range]
For $0<s<1$, determine the phase regularity and the correct small-phase
expansion. The periodized $H_s$ diverges, and the atomic fixed branch is
no longer described by the positive eigenfunction used here. The
appropriate normalization and dominant spectral component must be
identified before perturbing.
\end{question}

\begin{question}[Regularity at nonzero phases]
Classify phase regularity away from the atomic point. Does the answer
depend on the orbit of the zero of the weight, for example on its
periodicity or recurrence? The atomic result rules out a global analytic
theorem, but it does not rule out stronger regularity on substantial
subsets of the punctured phase circle.
\end{question}

\begin{question}[A complete fractional expansion]
Determine the higher-order structure of $E_{b,s}(c)$. Does the pressure
admit a convergent or asymptotic expansion involving even integer powers,
$|c|^s$, their products, and possible resonant logarithms? The exact
normal form isolates the unknown function and should be a better starting
point than differentiating the singular potential directly.
\end{question}

\begin{question}[Endpoint H\"older regularity]
For $k=\lceil s\rceil-1$ and non-even $s$, determine whether
$P_{b,s}\in C^{k,s-k}$ on a full neighborhood when $s$ is nonintegral,
and whether it belongs to $C^{s-1,1}$ when $s$ is odd. The derivative
asymptotic proves the pointwise order at zero, but a two-point endpoint
estimate needs additional control of $E_{b,s}$.
\end{question}

\begin{question}[A uniform critical crossover]
Develop a two-parameter asymptotic as $s\downarrow1$ and $c\to0$
simultaneously. The fixed-$s$ theorem is not uniform in that limit.
A useful result would identify the correct combined scaling variable and
match the supercritical cusp to the critical law.
\end{question}

\begin{question}[The full atomic spectrum]
The relaxation theorem determines the sharp norm-decay scale. Determine
the complete spectrum on $C^\alpha$ and on smoother or weighted spaces,
including the mechanism behind the logarithmic factor at
$s=1+\alpha$. Distinguish genuine eigenvalues, essential spectrum, and
endpoint norm effects rather than inferring one from the other.
\end{question}

\begin{question}[Other masks and multiple zeros]
Generalize the method to normalized trigonometric masks whose zeros have
different multiplicities, and to several forbidden inverse branches.
Which hypotheses yield a computable positive atomic eigenfunction, and
when is the first nonanalytic exponent determined by a zero multiplicity?
The arithmetic zeta sum in this paper may be replaced by a more general
lattice or orbit sum.
\end{question}

\begin{question}[Certified numerics and formalization]
Turn the explicit branch bounds into quantitative constants and combine
them with interval arithmetic to certify a usable phase neighborhood and
finite-size error. A proof-assistant development could then verify the
periodization, the sharp power-sum estimate, the spectral separation,
and the exact scalar identity in separate modules. This would convert
the ordinary-mathematics theorem into the repository's trusted setting.
\end{question}

\appendix
\section{A self-contained coarse Bernstein estimate}
\label{app:bernstein}
Here is the convolution argument used in \cref{lem:bernstein}. Let
\[
 F_M(x)=\frac1M\left(\frac{\sin\pi Mx}{\sin\pi x}\right)^2
\]
be the Fej\'er kernel. Its Fourier coefficients are
$(1-|k|/M)_+$, as follows by expanding the squared geometric sum. On
$|x|\le1/2$, direct differentiation gives
\begin{equation}
 |F_M'(x)|\le
 \begin{cases}
 CM^2,&|x|\le M^{-1},\\
 C\bigl(|x|^{-2}+M^{-1}|x|^{-3}\bigr),&M^{-1}<|x|\le1/2.
 \end{cases}
 \label{eq:Fejer-derivative}
\end{equation}
For the first bound one may differentiate the finite Fourier expansion:
$\sum_{|k|<M}|k|(1-|k|/M)=O(M^2)$. The second follows from the quotient
formula and $|\sin\pi x|\ge2|x|$ on this interval. Integrating shows
$\|F_M'\|_1\le CM$.

The de la Vall\'ee Poussin kernel $V_D=2F_{2D}-F_D$ has Fourier
coefficient $1$ for every $|k|\le D$. Hence $Q=Q*V_D$ for every
trigonometric polynomial of degree at most $D$, and
\[
 \|Q'\|_p=\|Q*V_D'\|_p
 \le\|Q\|_p\|V_D'\|_1\le CD\|Q\|_p\qquad(p\ge1)
\]
by the integral triangle inequality. Only a constant proportional to
$D$ is needed for pressure, so no sharp Bernstein constant is required.

\section{Proof dependencies and a formalization route}
\label{app:dependencies}
The central dependency chain is short:
\[
 \begin{gathered}
 \text{telescoping masks}
 \Longrightarrow\text{sharp relaxation}\\[2pt]
 \Longrightarrow\text{simple eigenvalue perturbation}
 \Longrightarrow\text{fixed-point normal form}.
 \end{gathered}
\]
The polynomial variation bound is a separate bridge identifying the
spectral eigenvalue with the cylinder definition of pressure. The zeta
sum evaluates the coefficient after the normal form is established.
The integer cancellation and the noninteger differentiability theorem
then follow from one-variable calculus.

A practical Lean integration could use the following boundaries.
First formalize the finite mask and iterate identities, keeping the
removable origin explicit. Next prove summability and the $C^1$ gluing
of $H_s$. Then formalize the invariant ideal estimate with an explicit
finite power sum; the lower bound uses only nonnegative summands of a
single test function. The spectral step should state and prove a general
isolated-rank-one perturbation lemma on a complex Banach space, rather
than admitting a specialized Perron theorem. Finally, package the scalar
identity at zero and the pressure comparison as independent interfaces.

None of these modules is supplied as already checked code. The present
package is designed to expose the theorem boundaries and the delicate
points: H\"older translation at a strict exponent, uniform positivity of
the eigenfunction, pressure identification in the presence of zeros, and
two-sided nonexistence of the next derivative.

\section{Source and convention ledger}
\label{app:sources}
The repository was inspected at commit
\nolinkurl{afb2d1227d8bc5df3beffc1463e958db3544960b} on 29 September
2026. The main predecessor is located under
\begin{quote}\small\raggedright
\nolinkurl{Analysis/FabiusFunction/docs/semi-formalized-research-frontiers/drafts/thue-morse/Thue_Morse_Integer_Pressure/}
\end{quote}
The primary-source check used arXiv:2509.22109v1 and
arXiv:2603.19001v1. The latter's page~4, including the stronger-regularity
question, was inspected in the rendered PDF. A separate machine-readable
ledger records URLs, scope, and the snapshot.

The essential convention conversions are
\[
 s=2q,\qquad p_q(c)=P_{2,2q}(c),\qquad
 a_{\mathrm{sine}}=c_{\mathrm{cosine}}+\tfrac12\pmod1.
\]
The transfer operator is an unnormalized sum, and the pressure uses
natural logarithms. The finite moment growth includes the additional
term $(s-1)\log b$ in \cref{eq:moment-asymp}. These conventions prevent
both the half-period ambiguity and the missing factor of $\log b$.

\begin{thebibliography}{9}
\bibitem{GKS}
P.~Gohlke, M.~Kesseb\"ohmer, and T.~I.~Schindler,
\emph{Generalized Thue--Morse measures: spectral and fractal analysis},
2025, \href{https://arxiv.org/abs/2509.22109v1}{arXiv:2509.22109v1}.

\bibitem{GLS}
P.~Gohlke, G.~Lamprinakis, and J.~Schmeling,
\emph{On a family of singular potentials: Parameter dependence of
thermodynamic characteristics}, 2026,
\href{https://arxiv.org/abs/2603.19001v1}{arXiv:2603.19001v1}.

\bibitem{Repo}
V.~Reshetnikov, \emph{ProveIt}, repository snapshot
\code{afb2d1227d8bc5df3beffc1463e958db3544960b}, inspected
29 September 2026.
\href{https://github.com/VladimirReshetnikov/ProveIt/tree/afb2d1227d8bc5df3beffc1463e958db3544960b}{Repository snapshot}.

\bibitem{IntegerRepo}
ProveIt research corpus,
\emph{Integer Pressure and a Missing Taylor Coefficient:
Analytic phase dependence and the sixth-moment optimum for generalized
Thue--Morse products}, research draft, 28 September 2026.
\href{https://github.com/VladimirReshetnikov/ProveIt/blob/afb2d1227d8bc5df3beffc1463e958db3544960b/Analysis/FabiusFunction/docs/semi-formalized-research-frontiers/drafts/thue-morse/Thue_Morse_Integer_Pressure/article.tex}{Pinned LaTeX source}.
\end{thebibliography}
\end{document}
